\documentclass[a4paper,11pt]{report}
\usepackage{a4wide}
\usepackage[frenchb]{babel}
\usepackage[utf8]{inputenc}
\usepackage[T1]{fontenc}
\usepackage{geometry}
\usepackage{graphicx}
\usepackage{fancybox}
\usepackage[pdftex]{hyperref} %Liens dans table de matieres
\usepackage{fancyhdr}
\usepackage{geometry}
\usepackage{url}
\usepackage{multicol}

\usepackage{pstricks}

%\geometry{scale=0.75, nohead}


\newcommand*{\enConstruction}[1]{
	\vspace*{.1cm}
	\doublebox{
	    \begin{minipage}{.9\linewidth}
	      --- #1 \\ 
	    \end{minipage}
	}
}

\newcommand*{\precision}[1]{
	\vspace*{.1cm}
	\fbox{
		\begin{minipage}{.9\linewidth}
			#1 
		\end{minipage}
  	}\\
}

% Title Page
\title{
	\includegraphics[scale=.2]{logofds.eps}\\
	\vspace*{1cm}
	TER FMIN200 \\ 
	-- \\
	Développement d'un jeu de type Bomberman en réseau sous Android et iOS
}

\author{BONVILA Olivier \and COUSEIN Kilian \and PITIOT Ludovic \and TARDIEU Benjamin}

\date{}

\begin{document}
\maketitle

\thanks{

  Nous tenons à remercier M. Giroudeau qui a toléré notre fâcheuse
  tendance à arriver en retard dans ses cours.
  
  Nous remercions encore une fois M. Giroudeau pour les points bonus
  qu'il va nous rajouter sur ce projet en échange de tous les journaux
  que nous lui apportions les jeudis matin.
  
  Enfin, nous remercions nos parents, frères, soeurs, proches, et plus
  particulièrement le monsieur avec les lunettes bizarres de la
  pizzeria du RU.\\ 

}


\tableofcontents

\chapter{Introduction}


\chapter{Présentation de l'application}

	\section{Moteurs}
		
		\subsection{Moteur de rendu}
		
			\subsubsection{Principes}
		
				\paragraph{Tiles}
			
				\paragraph{Sprites}
		
			\subsubsection{Objets}
			
			\subsubsection{Personnages}
			
		\subsection{Moteur physique}
		
			\subsubsection{Description}
		
			\subsubsection{Collisions}
			
			\subsubsection{Intéractions avec l'environnement}
			
		\subsection{Intelligence artificielle}

			\subsubsection{Méthode utilisée}
			
			\subsubsection{Facile}
			
			\subsubsection{Moyenne}
			
			\subsubsection{Difficile}

\chapter{Problèmes et solutions}

	Portage 
		 Opengl
		 Sérialisation/iOS
		 Adapter à toutes les résolutions d'écrans

\chapter{Résultats}


\end{document}      
